\documentclass[a4paper]{article}
\usepackage[autostyle]{csquotes}
\usepackage{fix-cm}
\usepackage[T2A]{fontenc}
\usepackage[utf8]{inputenc}
\usepackage[russian]{babel}
\usepackage{datetime}
\usepackage{mathtools}
\usepackage{marginnote}
\usepackage{geometry}
\usepackage{titlesec}
\usepackage{nicefrac}
\usepackage{bm}
\usepackage{listings}
\usepackage{bashful}
\usepackage{pgffor}
\usepackage{amsmath}
\usepackage{amsthm}
\usepackage{amsfonts}
\usepackage{enumitem}
\usepackage{svg}
\usepackage{color}
\usepackage{xcolor}
\usepackage{colortbl}
\usepackage{tabularx}
\usepackage{multirow}
\usepackage{multicol}
\usepackage{makecell}
\usepackage{booktabs}
\usepackage{array}
\usepackage[normalem]{ulem}
\usepackage{environ}
% \usepackage[noae]{Sweave}

\geometry{
    a4paper,
    top=20mm,
    right=15mm,
    bottom=20mm,
    left=30mm,
}

\setcounter{secnumdepth}{4}

\titleformat{\paragraph}
{\normalfont\normalsize\bfseries}{\theparagraph}{1em}{}
\titlespacing*{\paragraph}
{0pt}{3.25ex plus 1ex minus .2ex}{1.5ex plus .2ex}

\delimitershortfall-1sp

\newcommand\abs[1]{\left|#1\right|}
\newcommand\norm[1]{\left|\left|#1\right|\right|}
\newcommand\cond[1]{\text{cond}\left(#1\right)}
\newcommand{\m}[1]{\ensuremath{\bm{#1}}}
\newcommand\eqdef{\mathrel{\stackrel{\makebox[0pt]{\mbox{\normalfont\tiny def}}}{=}}}

\newcommand\numberthis{\addtocounter{equation}{1}\tag{\theequation}}

\newenvironment{where}[1][где:]
    {#1 \begin{tabular}[t]{>{}<{} @{${}{}$}l}}
    {\end{tabular}\\[\belowdisplayskip]}

\lstset{%
    basicstyle=\ttfamily,
    keywordstyle=\color{blue}\ttfamily,
    stringstyle=\color{red}\ttfamily,
    commentstyle=\color{green}\ttfamily,
    morecomment=[l][\color{magenta}]{\#},
    showstringspaces=false,
}

\DeclareTextFontCommand{\emph}{\boldmath\bfseries}

\begin{document}

% \input{main-concordance}
\begin{titlepage}
    \begin{center}
        \setlength{\parindent}{0cm}
        \fontsize{16pt}{16pt}\selectfont

        \textbf{ИНСТИТУТ ТРАНСПОРТА И СВЯЗИ}

        \rule{\textwidth}{1pt}

        \vspace*{2.0cm}

        \includegraphics[width=5cm]{tsi_logo}

        \vspace*{2.0cm}

        ФАКУЛЬТЕТ КОМПЬЮТЕРНЫХ НАУК И ТЕЛЕКОММУНИКАЦИЙ

        \vspace*{2.0cm}

        Лабораторная работа №1 \\
        По дисциплине \\
        <<Численные методы и прикладное программирование>>

        \vspace*{2.0cm}
        \fontsize{14pt}{14pt}\selectfont
        Тема: \\
        <<Методы решения системы линейных уравнений. \\ Число обусловленности матрицы>>

        \vfill
        \fontsize{12pt}{12pt}\selectfont

        \begin{flushright}
            Работу выполнили: \\
            \vspace*{0.25cm}
            {\itshape Дзенис Ричард} \\
            {\itshape Кобелев Денис} \\
            {\itshape Якушин Владислав} \\
        \end{flushright}

        \vfill
        
        Рига \\ 2017 г.
    \end{center}
\end{titlepage}

\tableofcontents

\pagebreak

\section{Формулировка задания}
\begin{itemize}
    \item Реализовать программным путём метод исключения Гаусса и итерационный метод Гаусса-Зейделя. 
    \item Результат работы программы проверить с помощью предоставленных примеров.
    \item Ручным или программным путём рассчитать число обусловленности матриц для примеров \eqref{slau:3} и \eqref{slau:5}.
    \item Для расчёта обусловленности выбрать Манхэттенскую или Евклидову норму.
    \item Составить отчёт с результатами вычислений и выводами, содержащими сравнение двух реализованных методов, а так же объяснить значения полученный при вычислении числа обусловленности матриц.
\end{itemize}

\subsection{Примеры}

\begin{equation}
    \begin{cases*}
           x_1  -2 x_2 + x_3 = 2 \\
         2 x_1  -5 x_2 - x_3 = -1 \\
        -7 x_1  +  x_3       = -2
    \end{cases*} \label{slau:1}
\end{equation}

\begin{equation}
    \begin{cases}
         5 x_1 -5 x_2 -3 x_3 +4 x_4 = -11 \\
           x_1 -4 x_2 +6 x_3 -4 x_4 = -10 \\
        -2 x_1 -5 x_2 +4 x_3 -5 x_4 = -12 \\
        -3 x_1 -3 x_2 +5 x_3 -5 x_4 = 8
    \end{cases} \label{slau:2}
\end{equation}

\begin{equation}
    \begin{cases*}
        2 x_1 -   x_2 -   x_3 = 5 \\
          x_1 + 3 x_2 - 2 x_3 = 7 \\
          x_1 + 2 x_2 + 3 x_3 = 10
    \end{cases*} \label{slau:3}
\end{equation}

\begin{equation}
    \begin{cases}
         8 x_1 + 5 x_2  + 3  x_3 = 30 \\
        -2 x_1 + 8 x_2  +    x_3 = 15 \\
           x_1 + 3 x_2  - 10 x_3 = 42
    \end{cases} \label{slau:4}
\end{equation}

\begin{equation}
    \begin{cases}
        0.78  x_1 + 0.563 x_2 = 0.217 \\
        0.913 x_1 + 0.659 x_2 = 0.254
    \end{cases} \label{slau:5}
\end{equation}

\section{Метод исключения Гаусса с ведущим элементом}

\subsection{Листинг}
\lstinputlisting[language=c++, firstline=5, lastline=49]{../gauss.cpp}

\subsection{Результаты работы алгоритма}

\subsubsection{Результат работы алгоритма на примере \eqref{slau:1}}
\bash[stdout]
../solve_problem gauss 1
\END

\subsubsection{Результат работы алгоритма на примере \eqref{slau:2}}
\bash[stdout]
../solve_problem gauss 2
\END

\subsubsection{Результат работы алгоритма на примере \eqref{slau:3}}
\bash[stdout]
../solve_problem gauss 3
\END

\subsubsection{Результат работы алгоритма на примере \eqref{slau:4}}
\bash[stdout]
../solve_problem gauss 4
\END

\subsubsection{Результат работы алгоритма на примере \eqref{slau:5}}
\bash[stdout]
../solve_problem gauss 5
\END

\section{Метод Гаусса-Зейделя}

\subsection{Листинг}
\lstinputlisting[language=c++, firstline=5, lastline=35]{../gauss-seidel.cpp}

\subsection{Результаты работы алгоритма с точностью $10^{-3}$}

\subsubsection{Результат работы алгоритма на примере \eqref{slau:1}}
\bash[stdout]
../solve_problem gauss-seidel 1 0.001
\END
\newline

\subsubsection{Результат работы алгоритма на примере \eqref{slau:2}}
\bash[stdout]
../solve_problem gauss-seidel 2 0.001
\END

В первых двух примерах не выполняется условие сходимости для итерационного метода, что и было успешно обнаружено программой.

\subsubsection{Результат работы алгоритма на примере \eqref{slau:3}}
\bash[stdout]
../solve_problem gauss-seidel 3 0.001
\END

\subsubsection{Результат работы алгоритма на примере \eqref{slau:4}}
\bash[stdout]
../solve_problem gauss-seidel 4 0.001
\END

\subsubsection{Результат работы алгоритма на примере \eqref{slau:5}}
\bash[stdout]
../solve_problem gauss-seidel 5 0.001
\END

\section{Определение невязки}
\lstinputlisting[language=python, firstline=32, lastline=34]{../find_residual.py}

\subsection{Полученные результаты для метода исключения Гаусса}

\paragraph{Пример \eqref{slau:1}}
\bash[stdout]
../find_residual.py 1 gauss
\END

\paragraph{Пример \eqref{slau:2}}
\bash[stdout]
../find_residual.py 2 gauss
\END

\paragraph{Пример \eqref{slau:3}}
\bash[stdout]
../find_residual.py 3 gauss
\END

\paragraph{Пример \eqref{slau:4}}
\bash[stdout]
../find_residual.py 4 gauss
\END

\paragraph{Пример \eqref{slau:5}}
\bash[stdout]
../find_residual.py 5 gauss
\END

\subsection{Полученные результаты для метода Гаусса-Зейделя}

\paragraph{Пример \eqref{slau:3}}
\bash[stdout]
../find_residual.py 3 gauss-seidel 0.001
\END

\paragraph{Пример \eqref{slau:4}}
\bash[stdout]
../find_residual.py 4 gauss-seidel 0.001
\END

\paragraph{Пример \eqref{slau:5}}
\bash[stdout]
../find_residual.py 5 gauss-seidel 0.001
\END

\section{Экспериментальное определение числа обусловленности матрицы}

\subsection{Часть листинга для нахождения числа обусловленности}
\lstinputlisting[language=python, firstline=35, lastline=49]{../find_cond.py}

\subsection{Полученные результаты}

При нахождении числа обусловленности экспериментальным методом, для расчётов неизвестных ($\m{X}$) использовался метод исключения Гаусса, для получения более высокой точности. А также, из-за того, что итерационным методом некоторые СЛАУ не возможно решить.

\paragraph{Пример \eqref{slau:1}}
\bash[stdout]
../find_cond.py 1 gauss
\END

\paragraph{Пример \eqref{slau:2}}
\bash[stdout]
../find_cond.py 2 gauss
\END

\paragraph{Пример \eqref{slau:3}}
\bash[stdout]
../find_cond.py 3 gauss
\END

\paragraph{Пример \eqref{slau:4}}
\bash[stdout]
../find_cond.py 4 gauss
\END

\paragraph{Пример \eqref{slau:5}}
\bash[stdout]
../find_cond.py 5 gauss
\END

\section{Дополнительные испытания}
В ходе сдачи лабораторной работы было необходимо провести дополнительные испытания на следующей матрице:
\begin{align}
    \begin{cases}
        x_{11}+ 2x_{12}  +3x_{13} = 10 \\
        x_{21}+ 3x_{22}  -2x_{23} = 7 \\
        2x_{31} -x_{32}  +x_{33} = 5
    \end{cases}
\end{align}

\paragraph{Решение СЛАУ с помощью метода исключения Гаусса}
\bash[stdout]
../solve_problem gauss test
\END

\paragraph{Решение СЛАУ с помощью метода Гаусса-Зейделя}
\bash[stdout]
../solve_problem gauss-seidel test 0.001
\END

Как можно видеть из результатов, данный алгоритм не сходится. Это обусловлено тем, что для данной СЛАУ не выполняется необходимое условие сходимости.

При преобразовании СЛАУ, путем перестановки третьего и первого уравнения местами, начинает выполнятся условие сходимости:
\begin{align}
    \begin{cases}
        2x_{31} -x_{32}  +x_{33} = 5 \\
        x_{11}+ 2x_{12}  +3x_{13} = 10 \\
        x_{21}+ 3x_{22}  -2x_{23} = 7
    \end{cases}
\end{align}

\bash[stdout]
../solve_problem gauss-seidel test2 0.001
\END

\section{Выводы}
В ходе выполнения данной лабораторной работы были реализованы два алгоритма нахождения решения СЛАУ на языке C++:
\begin{itemize}
    \item Прямой метод: <<метод исключения Гаусса>>;
    \item Итерационный метод: <<метод Гаусса-Зейделя>>.
\end{itemize}
А также были вычислены числа обусловленности для заданных примеров экспериментальным путём.

\bigskip

При сравнении двух реализованных методов можно заметить, что метод Гаусса-Зейделя не предназначен для решения всевозможных СЛАУ.
При проверки работы алгоритмов на предоставленных примерах, в 1-ом и 2-ом примере метод Гаусса-Зейделя не сходиться. Примеры 3 и 4 имели довольно схожие результаты.
5-ый пример при низкой точности ($\varepsilon \approx 10^{-6}$), результаты, полученные методом Гаусса-Зейделя довольно сильно расходиться с методом исключения Гаусса. 

\bigskip

Разница времени выполнения обоих методов для данных (малых) СЛАУ была довольно незначительной, поэтому, для сравнения времени выполнения, было принято решение создать СЛАУ размером $1500 \times 1500$, и произвести замеры на ней.

Каждый алгоритм был запущен $10$ раз на одной и той-же матрице $1500 \times 1500$, в результате получено:
\begin{center}
    \begin{tabular}{r|cc}
        Метод & Среднее время выполнения (с) & Точность \\
        \hline
        Метод исключения Гаусса & $5.357956087$ & $\pm 0.33\%$ \\
        Метод Гаусса-Зейделя    & $0.655645297$ & $\pm 0.73\%$ \\
    \end{tabular}
\end{center}

Полученные результаты были ожидаемы, т.к. метод исключения Гаусса имеет алгоритмическую сложность $O(n^3)$, тогда как метод Гаусса-Зейделя: $O(mn^2)$, где $m$ -- количество итераций, которое необходимо для достижения требуемой точности.
\begin{equation*}
    O(n^3) > O(mn^2) \qquad \text{если} \; n \gg m
\end{equation*}

Таким образом при не больших СЛАУ (когда $n \approx m$) метод исключения Гаусса не медленнее метода Гаусса-Зейделя. Однако при увеличении размера СЛАУ, когда $n \gg m$, первый метод начинает работать значительно медленнее второго.

\bigskip

В ходе вычисления числа обусловленности Матриц была использована Евклидова норма. За исключением 5-ого примера, все числа обусловленности получились до 10, что означает хорошую обусловленность.
Число обусловленности матрицы из 5-ого уравнения составила $\approx 227240$, что явно свидетельствует о плохой обусловленности, так как значение выше 1000. Скорее всего, данный факт и является причиной расхождения результатов метода Гаусса-Зейделя и метода исключений Гаусса при низкой точности.

\end{document}

